% BHU PYQ -- Transcribed & Corrected
% Paper: ENGAE-31: Reading Skills (Mid Term)
% Exam: B.A. / B.Sc. / B.Com. (Hons.) Semester III (Mid Term) Examination, 2025-26 (NEP)
% Category: Ability Enhancement Course (AEC)

\documentclass[11pt,a4paper]{article}
\usepackage[a4paper,top=2.5cm,bottom=2.5cm,left=2.8cm,right=2.8cm,headheight=14pt]{geometry}
\usepackage{amsmath,amssymb}
\usepackage[T1]{fontenc}
\usepackage[utf8]{inputenc}
\usepackage{lmodern,microtype,fancyhdr,enumitem,hyperref,lastpage,parskip}

\hypersetup{
    colorlinks=true,
    urlcolor=black,
    linkcolor=black,
    pdftitle={ENGAE-31: Reading Skills (Mid Term) -- BHU Sem III 2025-26 (NEP)}
}

\pagestyle{fancy}
\fancyhf{}
\renewcommand{\headrulewidth}{0.4pt}
\renewcommand{\footrulewidth}{0.4pt}
\fancyhead[L]{\small \textit{Banaras Hindu University}}
\fancyhead[R]{\small \textit{ENGAE-31: Reading Skills (Mid Term)}}
\fancyfoot[C]{\small Page \thepage\ of \pageref{LastPage}}


\newlist{parts}{enumerate}{2}
\setlist[parts,1]{label=(\Alph*),leftmargin=*,itemsep=3pt,topsep=2pt}

\newcommand{\pts}[1]{\hfill \textit{[#1]}}


\begin{document}


\begin{center}
  {\large\bfseries BANARAS HINDU UNIVERSITY}\\[2pt]
  {\normalsize B.A. / B.Sc. / B.Com. (Hons.) --- Semester III Mid Term Examination, 2025-26 (NEP)}\\[6pt]
  \rule{0.8\linewidth}{0.5pt}\\[6pt]
  {\large\bfseries Ability Enhancement Course (AEC)}\\[4pt]
  {\large\bfseries Paper Code: ENGAE-31 : Reading Skills (Mid Term)}\\[4pt]
  {\normalsize Time: 1 Hour\hfill Maximum Marks: 30}
\end{center}

\rule{\linewidth}{0.4pt}

\smallskip

\noindent \textit{Instructions: Answer all 30 objective questions. Each question carries equal marks (1 mark each).}

\bigskip




\noindent \textbf{Question 1.} Which of the following is a primary goal of assessing your reading level?\pts{1}

\begin{parts}
  \item To compare your reading speed with others
  \item To understand your current comprehension ability
  \item To memorize new vocabulary
  \item To analyse the author's writing style
\end{parts}

\medskip




\noindent \textbf{Question 2.} Which of the following is NOT a common barrier to reading comprehension?\pts{1}

\begin{parts}
  \item Limited vocabulary
  \item Poor concentration
  \item Reading with purpose
  \item Lack of background knowledge
\end{parts}

\medskip




\noindent \textbf{Question 3.} Which kind of reading is mainly involved in understanding a pie chart or a bar graph?\pts{1}

\begin{parts}
  \item Visual
  \item Verbal
  \item Symbolic
  \item Narrative
\end{parts}

\medskip




\noindent \textbf{Question 4.} The purpose of finding errors in a passage is to:\pts{1}

\begin{parts}
  \item Improve your handwriting
  \item Enhance spelling and grammatical accuracy
  \item Memorize a particular text
  \item Increase your reading speed
\end{parts}

\medskip




\noindent \textbf{Question 5.} What is the implied meaning of the phrase ``He has a heart of stone''?\pts{1}

\begin{parts}
  \item He has a medical problem
  \item He loves stones
  \item He is strong and healthy
  \item He is unkind or unemotional
\end{parts}

\medskip




\noindent \textbf{Question 6.} To improve your reading comprehension skills, you can:\pts{1}

\begin{parts}
  \item Select a text slightly above your current comfort level
  \item Re-read only the easiest texts repeatedly
  \item Avoid challenging or unfamiliar topics
  \item Read without reflecting on meaning
\end{parts}

\medskip




\noindent \textbf{Question 7.} A friend of yours struggles with staying focused during reading sessions, due to a limited attention span. Which of the following barriers is she facing?\pts{1}

\begin{parts}
  \item Emotional barrier
  \item Cognitive barrier
  \item Physical barrier
  \item Social barrier
\end{parts}

\medskip




\noindent \textbf{Question 8.} Your friend recently received some upsetting news from home. Now, he is finding it difficult to read in class. Which of the following barriers is he facing?\pts{1}

\begin{parts}
  \item Emotional barrier
  \item Cognitive barrier
  \item Physical barrier
  \item Social barrier
\end{parts}

\medskip




\noindent \textbf{Question 9.} The word ``recieve'' instead of ``receive'' is an example of:\pts{1}

\begin{parts}
  \item Grammatical error
  \item American and British conventions
  \item Spelling error
  \item Semantic confusion
\end{parts}

\medskip




\noindent \textbf{Question 10.} When a story says ``The sun smiled on the children,'' the sentence means:\pts{1}

\begin{parts}
  \item The sun was literally smiling
  \item The weather was pleasant
  \item The author made a mistake
  \item The children were smiling
\end{parts}

\medskip




\noindent \textbf{Question 11.} Reading aloud with appropriate speed, pauses, and tone helps improve:\pts{1}

\begin{parts}
  \item Writing fluency
  \item Listening comprehension only
  \item Silent reading skills
  \item Reading fluency and pronunciation
\end{parts}

\medskip




\noindent \textbf{Question 12.} Which of the following strategies helps overcome vocabulary-related barriers to reading?\pts{1}

\begin{parts}
  \item Ignoring difficult words
  \item Skipping the text
  \item Reading faster
  \item Using contextual clues and dictionary
\end{parts}

\medskip




\noindent \textbf{Question 13.} Which of the following kinds of texts requires a keen understanding of rhyme and rhythm:\pts{1}

\begin{parts}
  \item Poetry
  \item Essays
  \item Newspaper
  \item Book report
\end{parts}

\medskip




\noindent \textbf{Question 14.} The main purpose of choral reading is to:\pts{1}

\begin{parts}
  \item Compete individually
  \item Read silently
  \item Develop group rhythm and expression
  \item Memorize a poem
\end{parts}

\medskip




\noindent \textbf{Question 15.} The main aim of silent, deep reading is to:\pts{1}

\begin{parts}
  \item Improve articulation
  \item Read aloud at a faster pace
  \item Focus on comprehension and interpretation
  \item Perform before others
\end{parts}

\medskip




\noindent \textbf{Question 16.} Effective note-taking helps a reader:\pts{1}

\begin{parts}
  \item Memorise the text line by line
  \item Identify and retain key ideas
  \item Avoid understanding
  \item Rewrite the full text
\end{parts}

\medskip




\noindent \textbf{Question 17.} Skimming a newspaper article helps to:\pts{1}

\begin{parts}
  \item Memorise details
  \item Get the general idea quickly
  \item Identify every important name
  \item Analyse the literary style
\end{parts}

\medskip




\noindent \textbf{Question 18.} If you are quickly checking a menu for sugar-free options, you are practicing:\pts{1}

\begin{parts}
  \item Deep reading
  \item Skimming
  \item Scanning
  \item Predictive reading
\end{parts}

\medskip




\noindent \textbf{Question 19.} Intensive reading focuses on:\pts{1}

\begin{parts}
  \item Understanding a particular text in detail
  \item Reading mainly for pleasure
  \item Reading only the key points
  \item Engaging with the pictorial illustrations
\end{parts}

\medskip




\noindent \textbf{Question 20.} Reading many articles on a topic, instead of focusing on (and analysing) one short article, is an example of:\pts{1}

\begin{parts}
  \item Intensive reading
  \item Extensive reading
  \item Predictive reading
  \item Skimming
\end{parts}

\medskip




\noindent \textbf{Question 21.} The process of learning new words while reading is called:\pts{1}

\begin{parts}
  \item Vocabulary reinforcement
  \item Word mapping
  \item Contextual vocabulary building
  \item Mnemonics
\end{parts}

\medskip




\noindent \textbf{Question 22.} The process of Predictive Reading involves:\pts{1}

\begin{parts}
  \item Skipping difficult words for now
  \item Reading the summary first
  \item Using the dictionary for each line
  \item Guessing the meaning of new words from context
\end{parts}

\medskip




\noindent \textbf{Question 23.} Using the sentence ``My Very Excellent Mate Just Served Us Nachos'' to remember the planets Mercury, Venus, Earth, Mars, Jupiter, Saturn, Uranus, and Neptune, is an example of:\pts{1}

\begin{parts}
  \item Mnemonic, memory aid
  \item A poetic rhyme
  \item A visual chart
  \item Note-taking method
\end{parts}

\medskip




\noindent \textbf{Question 24.} When a reader connects prior knowledge with clues from the text, they are engaging in:\pts{1}

\begin{parts}
  \item Summarisation
  \item Inferencing
  \item Paraphrasing
  \item Predicting
\end{parts}

\medskip




\noindent \textbf{Question 25.} Reflective reading primarily encourages:\pts{1}

\begin{parts}
  \item Memory retention
  \item Skimming and scanning
  \item Critical self-assessment and insight
  \item Quick recall of key points
\end{parts}

\medskip




\noindent \textbf{Question 26.} After reading an article on environmental pollution, a student relates it to their city's air quality. This is an example of:\pts{1}

\begin{parts}
  \item Reflective thinking
  \item Mnemonic
  \item Rhyme recognition
  \item Literary comprehension
\end{parts}

\medskip




\noindent \textbf{Question 27.} Personalising a text means:\pts{1}

\begin{parts}
  \item Editing it in your own style
  \item Relating it to your own experience
  \item Memorising every detail
  \item Skipping unknown parts
\end{parts}

\medskip




\noindent \textbf{Question 28.} A reader who imagines the physical setting of a story while reading is using:\pts{1}

\begin{parts}
  \item Visualisation
  \item Paraphrasing
  \item Memorisation
  \item Symbolism
\end{parts}

\medskip




\noindent \textbf{Question 29.} A mind map or flowchart used to represent main ideas from an essay is an example of:\pts{1}

\begin{parts}
  \item Graphic organisation
  \item Mnemonics
  \item Predictive reading
  \item Phonetic chart
\end{parts}

\medskip




\noindent \textbf{Question 30.} Which of the following can NOT be used to summarise a chapter:\pts{1}

\begin{parts}
  \item A few sentences on the main points
  \item A table with the main data
  \item A flow chart with key concepts
  \item A news clipping on the same topic
\end{parts}

\bigskip

\begin{center}
  \textbf{*** End of Question Paper ***}
\end{center}

\end{document}
